\section{Introduction}
\label{sect:intro}

Binary-star orbits are among the principal routes to dynamical stellar masses. Independent resolved relative astrometry constrains the sky-plane geometry, while a double-lined spectroscopic binary (SB2) supplies both line-of-sight velocity amplitudes and therefore strong information on the mass ratio. Gaia adds precise one-dimensional astrometric information over a mission-length time baseline \citep{GaiaMission2016}. The scientific issue addressed here is not how to combine those data types in general, but how the \emph{measurement model} used for a luminous close pair propagates into the inferred physical orbit and component masses.

The broad combined-orbit problem is already well developed. BINARYS jointly treats Hipparcos/Gaia absolute astrometry, independent relative astrometry, and component-labelled radial velocities \citep{Leclerc2023}, while orvara performs posterior inference from radial velocity, absolute astrometry, and relative astrometry \citep{Brandt2021}. \citet{Chevalier2023} used BINARYS to combine Gaia DR3 non-single-star solutions with SB2 data and derive individual masses and luminosities, demonstrating directly that Gaia astrometry plus SB2 spectroscopy is not itself a gap. Component-mass inference from unresolved Gaia astrometric orbits and photometric information is also established \citep{BailerJones2026}.

The unresolved-to-resolved transition remains more problematic. Gaia's published DR3 astrometric-binary processing used a preliminary treatment to discard partially resolved double stars before orbital modelling \citep{Halbwachs2023}, and \citet{Holl2023} demonstrated that close source structure can introduce scan-angle-dependent image-parameter biases and spurious signals. BINARYS is particularly informative in this context: it already models non-trivial Hipparcos transit response when the secondary contributes significant light, but its Gaia treatment does not extend the same response modelling to luminous partially resolved Gaia pairs; \citet{Leclerc2023} explicitly identify the need for more detailed Gaia information to treat that flux-contamination regime.

A nonlinear, separation-dependent Gaia-like blended-source response is also prior art. \citet{ElBadry2024} released \texttt{gaiamock}, whose \texttt{al\_bias\_binary} places the measured one-dimensional coordinate at the peak of the combined along-scan light profile. A deeper audit of the public source also finds a forward routine that predicts this response and a primary-star RV curve for the same stellar binary. We therefore make no claim that resolution-aware stellar Gaia modelling, or even response-aware Gaia-plus-RV forward prediction, originates here. Our equal-width implementation reproduces the same response family over their common single-peak domain.

Nor is response-aware orbital inference itself new in a broader astronomical sense. \citet{Liu2024} fitted the binary asteroid (4337) Arecibo using Gaia Focused Product Release astrometry with an unresolved/partially-resolved/resolved response and external occultation constraints, deriving the relative orbit and physical properties. Their analysis is sequential and is not a stellar SB2 joint posterior, but it removes any defensible claim that putting a partially resolved Gaia response into an orbital MCMC is new in general.

The response theory is also more general than the baseline surrogate used below. \citet{Penoyre2026} derived the position and resolvability of blended Gaussian point sources for an elongated Gaia-like PSF, reducing the problem to a one-dimensional blend with an orientation-dependent effective width. The published correction to that paper changes several signs, so any implementation must use the corrected equations. The constant-width equal-profile response in our frozen experiments is consequently a controlled restricted case rather than a physical Gaia PSF/LSF model or new resolvability theory.

Finally, Bayesian Gaia epoch-astrometry plus RV inference is already implemented in \texttt{kima} \citep{Baycroft2026}. Public Octofitter code likewise provides Gaia DR4 epoch-astrometry likelihood infrastructure alongside relative-astrometry and radial-velocity modules. In the specific public Gaia likelihood paths audited for this work, we did not identify the physical marginal-resolution blended-profile operator sought here; however, these software audits are evidence only about the inspected implementations, not proof of absence elsewhere.

Against that background, we investigate a deliberately narrow \emph{candidate inference gap}. Based on the published literature and publicly available implementations examined here, we did not identify a stellar target-level analysis in which a physical marginal-resolution Gaia blended-source response is placed directly inside an epoch-astrometric likelihood and inferred simultaneously with independent resolved relative astrometry and \emph{both} SB2 velocity curves. We do not claim priority for that intersection. The literature and source audit is recorded in \texttt{docs/literature\_gap.md} and should be repeated before submission.

More importantly, the scientific question does not depend on winning a checklist of software capabilities. We ask: \emph{what biases in individual stellar dynamical masses and parallax arise when marginal-resolution Gaia measurements are treated as an ordinary photocentre, and how accurate must the astrometric response model be before those biases and posterior-coverage failures disappear?} This response-fidelity question remains meaningful even if another framework later implements the same combination of data types.

The immediate goal is not a calibrated Gaia image model. Gaia's real PSF/LSF is time-, colour-, focal-plane-, and instrument-state dependent \citep{Rowell2021,Rowell2026}. We therefore retain the equal-width response as a transparent baseline and use it to establish the validation chain: physical orbit conventions, the unresolved photocentre limit, the single/multi-peak boundary of the restricted surrogate, mission-grounded scan schedules, a full-sky transition experiment, and a matched-transit-count control. The next stage is explicitly a response-misspecification hierarchy in which more general injected profiles are fitted with successively simplified measurement models and judged primarily by mass/parallax bias and posterior coverage rather than by $\Delta\chi^2$ alone.

The frozen 23\,000-fit and 11\,040-fit experiments remain baseline equal-width synthetic tests. They are not empirical Gaia calibration results and are not retroactively described as using the later absolute-astrometry, posterior-sampling, or response-misspecification extensions of the code.
